% lecturekit.tex: the study PDF's Beamer theme (Neural Models slide theme v0.3, as a Quarto include).
% Colours are the nm<Name> tokens of _brand.yml, written to _lecturekit/tokens.tex by the filter together with
% the fonts: TeX Gyre Heros (Helvetica) for sans and TeX Gyre Pagella (Palatino) for serif in the paper preset,
% loaded by file from the extension's fonts/, so a build is identical on every machine and needs no local font.
% Fall-back values below cover a deck rendered without a brand file.
\usepackage{amsmath}
\usepackage{fontspec}
\usepackage{textcase}
\usepackage{etoolbox}
\IfFileExists{unicode-math.sty}{\usepackage{unicode-math}}{}
\usepackage{tikz}
\usetikzlibrary{arrows.meta,positioning,calc,decorations.pathreplacing}
\usefonttheme{professionalfonts}

% ---------------------------------------------------------------- tokens from the brand, then the deck's metadata
\InputIfFileExists{_lecturekit/tokens.tex}{}{%
  % no token file: the engine's families by name, if the machine has them
  \IfFontExistsTF{TeX Gyre Pagella}{\setmainfont{TeX Gyre Pagella}[Ligatures=TeX]}{}%
  \IfFontExistsTF{TeX Gyre Heros}{\setsansfont{TeX Gyre Heros}[Ligatures=TeX]}{}%
  \IfFontExistsTF{TeX Gyre Pagella}{\newfontfamily\nmSerifCaps{TeX Gyre Pagella}[Ligatures=TeX,LetterSpace=2.0]}{}%
  \IfFontExistsTF{TeX Gyre Heros}{\newfontfamily\nmSansCaps{TeX Gyre Heros}[Ligatures=TeX,LetterSpace=8.0]}{}%
  \IfFontExistsTF{TeX Gyre Pagella Math}{\setmathfont{TeX Gyre Pagella Math}}{}}
\providecommand{\nmSerifCaps}{\rmfamily}
\providecommand{\nmSansCaps}{\sffamily}
\newif\iflkCitationSmall

\providecolor{nmInk}{HTML}{202B33}
\providecolor{nmMuted}{HTML}{566573}
\providecolor{nmHairline}{HTML}{C8C3BA}
\providecolor{nmPaper}{HTML}{F3F2EE}
\providecolor{nmSlide}{HTML}{FFFFFF}
\providecolor{nmInput}{HTML}{0072B2}
\providecolor{nmState}{HTML}{007F73}
\providecolor{nmUpdate}{HTML}{C45A00}
\providecolor{nmContext}{HTML}{7A5195}
% Teaching roles are separate from the model-variable colours in figures.
\providecolor{nmConcept}{HTML}{175B8E}
\providecolor{nmCaveat}{HTML}{A33A32}
\providecolor{nmExample}{HTML}{956000}
\providecolor{nmAnchor}{HTML}{006C65}
\providecolor{nmQuestion}{HTML}{70508C}
% The four pRF models of the workshop's advanced lecture (and the course's lectures 7 and 8); distinguishable under
% deutan, protan and tritan vision against each other and against nmInk (CIELAB difference at least 19.6).
\providecolor{nmGauss}{HTML}{175B8E}
\providecolor{nmDoG}{HTML}{B87400}
\providecolor{nmCSS}{HTML}{B63564}
\providecolor{nmDN}{HTML}{009999}
\providecolor{nmFuture}{HTML}{7A5195}
% optional tokens: heading (titles) and body (running text) fall back to ink
\providecolor{nmHeading}{named}{nmInk}
\providecolor{nmBody}{named}{nmInk}

% ---------------------------------------------------------------- colours and sizes
\setbeamercolor{normal text}{fg=nmBody,bg=nmSlide}
\setbeamercolor{structure}{fg=nmInput}
% The frame title and the canvas take their colours from two named colours that the role sets (nmRoleColor,
% nmRoleBg). The role is a frame option, parsed by keyval; keyval is not re-entrant, so the option's handler must
% not call \setbeamercolor (another \setkeys): with [fragile,nmrole=...] that ended in "fragile undefined".
\colorlet{nmRoleColor}{nmHeading}
\colorlet{nmRoleBg}{nmSlide}
\setbeamercolor{background canvas}{bg=nmRoleBg}
\setbeamercolor{frametitle}{fg=nmRoleColor,bg=nmRoleBg}
\setbeamercolor{title}{fg=nmHeading}
\setbeamercolor{alerted text}{fg=nmInk}
\setbeamerfont{title}{family=\rmfamily,size=\fontsize{26}{30}\selectfont,series=\mdseries}
\setbeamerfont{frametitle}{family=\rmfamily,size=\fontsize{20}{24}\selectfont,series=\mdseries}
\setbeamerfont{normal text}{size=\fontsize{12}{15}\selectfont}
\AtBeginDocument{\usebeamerfont{normal text}}
\setbeamersize{text margin left=9mm,text margin right=9mm}
\setbeamertemplate{navigation symbols}{}
\newcommand{\nmTitleCaps}[1]{{\nmSerifCaps\MakeTextUppercase{#1}}}
\newcommand{\nmLabelCaps}[1]{{\nmSansCaps\MakeTextUppercase{#1}}}
% Frame titles are uppercased from Beamer's raw title, with the "(cont.)" or roman continuation of
% allowframebreaks frames appended outside \MakeTextUppercase (uppercasing \insertframetitle fails on it).
\makeatletter
\newcommand{\nmFrameTitleCaps}{{\nmSerifCaps\MakeTextUppercase{\beamer@frametitle}%
  \ifnum\beamer@autobreakcount>0\relax\space\usebeamertemplate*{frametitle continuation}\fi}}
\makeatother
% Deck metadata beyond Beamer's own: a short running title for the foot and an eyebrow line above the main
% title. Markdown decks set these through the shorttitle and eyebrow keys (dual-output.lua writes
% _lecturekit/<deck>-deck.tex); LaTeX decks call them directly.
\newcommand{\nmShortTitleText}{}
\newcommand{\nmEyebrowText}{}
\newcommand{\nmShortTitle}[1]{\renewcommand{\nmShortTitleText}{#1}}
\newcommand{\nmEyebrow}[1]{\renewcommand{\nmEyebrowText}{#1}}
\newcommand{\nmFootTitle}{\ifx\nmShortTitleText\empty\inserttitle\else\nmShortTitleText\fi}
% Title page: eyebrow, capitals, a short rule, then subtitle, author and date; left aligned like every frame title.
\setbeamertemplate{title page}{%
  \vspace*{9mm}%
  \ifx\nmEyebrowText\empty\else
    {\sffamily\mdseries\fontsize{8.5}{10}\selectfont\color{nmMuted}\nmLabelCaps{\nmEyebrowText}\par}\vspace{4.5mm}%
  \fi
  {\usebeamerfont{title}\usebeamercolor[fg]{title}\nmTitleCaps{\inserttitle}\par}%
  \nointerlineskip\vspace{4.5mm}{\color{nmInk}\rule{14mm}{0.6pt}\par}%
  \nointerlineskip\vspace{4.5mm}{\rmfamily\mdseries\fontsize{15}{19}\selectfont\color{nmMuted}\insertsubtitle\par}%
  \vspace{8mm}{\sffamily\mdseries\fontsize{11}{14}\selectfont\color{nmBody}\insertauthor\par}%
  \ifx\insertinstitute\empty\else\vspace{1.2mm}{\sffamily\mdseries\fontsize{9}{11}\selectfont\color{nmMuted}\insertinstitute\par}\fi
  \vspace{1.2mm}{\sffamily\mdseries\fontsize{9}{11}\selectfont\color{nmMuted}\insertdate\par}}
\setbeamertemplate{itemize item}{\color{nmInput}\small\textbullet}
\setbeamertemplate{itemize subitem}{\color{nmMuted}\tiny\textbullet}
\setbeamerfont{itemize/enumerate body}{size=\fontsize{12}{15}\selectfont}
\setbeamerfont{itemize/enumerate subbody}{size=\fontsize{10.5}{13}\selectfont}

% ---------------------------------------------------------------- roles
\newcommand{\nmRole}{EXPLANATION}
\newcommand{\slidecue}[1]{\renewcommand{\nmRole}{#1}}
\newcommand{\nmRoleIcon}{}
\newcommand{\nmSetRole}[3]{%
  \slidecue{#1}\colorlet{nmRoleColor}{#2}\colorlet{nmRoleBg}{#3}\renewcommand{\nmRoleIcon}{}}
\pgfkeys{/nm/role/.is choice,
 /nm/role/explanation/.code={\nmSetRole{EXPLANATION}{nmHeading}{nmSlide}},
 /nm/role/concept/.code={\nmSetRole{CORE CONCEPT}{nmConcept}{nmSlide}},
 /nm/role/caveat/.code={\nmSetRole{CAVEAT / INTERPRET WITH CARE}{nmCaveat}{nmSlide}},
 /nm/role/example/.code={\nmSetRole{WORKED EXAMPLE / FOLLOW THE STEPS}{nmExample}{nmSlide}},
 /nm/role/goals/.code={\nmSetRole{LEARNING GOALS}{nmAnchor}{nmPaper}},
 /nm/role/summary/.code={\nmSetRole{TAKEAWAYS}{nmAnchor}{nmPaper}},
 /nm/role/question/.code={\nmSetRole{CHECK YOUR UNDERSTANDING}{nmQuestion}{nmSlide}},
 /nm/role/section/.code={\nmSetRole{NEXT PART}{nmAnchor}{nmPaper}},
 /nm/role/pause/.code={\nmSetRole{PAUSE THE RECORDING}{nmQuestion}{nmPaper}},
 /nm/role/detection/.code={\nmSetRole{MAPS}{nmConcept}{nmSlide}\renewcommand{\nmRoleIcon}{\nmIcon{nmConcept}{\draw[nmConcept,line width=0.45pt] plot[domain=-1.05:1.05,samples=31] (\x,{0.7*sin(\x*343)});}}},
 /nm/role/estimation/.code={\nmSetRole{ESTIMATION}{nmState}{nmSlide}\renewcommand{\nmRoleIcon}{\nmIcon{nmState}{\draw[nmState,line width=0.45pt] (0,0) circle (0.8); \fill[nmState] (0,0) circle (0.22);}}},
 /nm/role/discrimination/.code={\nmSetRole{DISCRIMINATION}{nmContext}{nmSlide}\renewcommand{\nmRoleIcon}{\nmIcon{nmContext}{\draw[nmContext,line width=0.45pt] (-1,-0.5) -- (0,0) -- (1,0.75) (0,0) -- (1,-0.6);}}},
 /nm/role/where/.code={\nmSetRole{WHERE}{nmHeading}{nmSlide}\renewcommand{\nmRoleIcon}{\nmIcon{nmHeading}{\draw[nmHeading,line width=0.4pt] (-1.1,0) -- (1.1,0) (0,-1.1) -- (0,1.1); \fill[nmHeading] (0.5,0.4) circle (0.33);}}},
 /nm/role/how/.code={\nmSetRole{HOW}{nmHeading}{nmSlide}\renewcommand{\nmRoleIcon}{\nmIcon{nmHeading}{\draw[nmHeading,line width=0.45pt] plot[domain=0:1,samples=21] ({2*\x-1},{1.5*sqrt(\x)-0.75});}}},
 /nm/role/what/.code={\nmSetRole{WHAT}{nmHeading}{nmSlide}\renewcommand{\nmRoleIcon}{\nmIcon{nmHeading}{\begin{scope}\clip (0,0) circle (1.2); \foreach \u in {-1.2,-0.4,0.4,1.2} \draw[nmHeading,line width=0.45pt] (\u-0.9,-1.4) -- (\u+0.9,1.4); \end{scope}}}},
 /nm/role/when/.code={\nmSetRole{WHEN}{nmHeading}{nmSlide}\renewcommand{\nmRoleIcon}{\nmIcon{nmHeading}{\draw[nmHeading,line width=0.45pt] (-1.1,-0.55) -- (-0.45,-0.55) -- (-0.45,0.6) -- (0.35,0.6) -- (0.35,-0.55) -- (1.1,-0.55);}}}}
% An objective's icon in the eyebrow: a small circle in the objective's colour with its glyph.
\newcommand{\nmIcon}[2]{\tikz[baseline=-0.62ex,x=0.9mm,y=0.9mm]{\draw[#1,line width=0.5pt,fill=#1!12] (0,0) circle (1.75); #2}\hspace{1.6mm}}
\newcommand{\sliderole}[1]{\pgfkeys{/nm/role/#1}}
% A frame-local option prevents one slide's role leaking into the next.
\makeatletter
\define@key{beamerframe}{nmrole}{\sliderole{#1}}
\makeatother
% Frame title block: role label, a short rule in the role colour, letterspaced capitals.
\setbeamertemplate{frametitle}{%
  \vspace{3.0mm}{\sffamily\mdseries\fontsize{9}{11}\selectfont\color{nmRoleColor}\nmRoleIcon\nmLabelCaps{\nmRole}\par}%
  \nointerlineskip\vspace{2.4mm}{\color{nmRoleColor}\rule{14mm}{0.6pt}\par}%
  \nointerlineskip\vspace{3.6mm}{\usebeamerfont{frametitle}\nmFrameTitleCaps\par}\vspace{1.0mm}}
% Running foot: hairline, short title, page of total. The title page keeps the same height but prints nothing.
\setbeamertemplate{footline}{%
  \hspace*{9mm}\begin{minipage}{\dimexpr\paperwidth-18mm\relax}\sffamily\mdseries\fontsize{7}{8}\selectfont
    \ifnum\value{framenumber}>1
      {\color{nmHairline}\rule{\linewidth}{0.3pt}}\par\nointerlineskip\vspace{1.2mm}%
      {\color{nmMuted}\makebox[\linewidth][s]{\nmLabelCaps{\nmFootTitle}\hfill\insertframenumber\,/\,\inserttotalframenumber}\strut\par}%
    \else
      {\rule{0pt}{0.3pt}}\par\nointerlineskip\vspace{1.2mm}{\strut\par}%
    \fi
  \end{minipage}\hspace*{9mm}\vspace{1.6mm}}
% Numbered points carry the role colour; body text stays charcoal.
\setbeamertemplate{enumerate item}{{\sffamily\mdseries\small\color{nmRoleColor}\ifnum\insertenumlabel<10 0\fi\insertenumlabel}}
\newcommand{\anchorline}[2]{%
  \par\vspace{2.6mm}\noindent
  \begin{minipage}[t]{.07\textwidth}\vspace{0pt}\color{nmRoleColor}\sffamily\small #1\end{minipage}%
  \begin{minipage}[t]{.91\textwidth}\vspace{0pt}\raggedright #2\end{minipage}\par}
\newcommand{\takehome}[1]{{\rmfamily\fontsize{18}{22}\selectfont #1\par}\vspace{3mm}}

% ---------------------------------------------------------------- figures
% TikZ figure files beside a deck: \nmfigure{assets/figures/name.tex} scales a figure down to the text width
% when it is wider (never up); \nmfigureh{4cm}{...} also caps the height. The filter turns both into an SVG for
% the browser. Implemented with \resizebox rather than adjustbox: Pandoc's template sets default includegraphics
% keys that adjustbox would apply to the figure as well.
\newsavebox{\nmFigBox}
\makeatletter
\newcommand{\nmFitBox}[1]{%
  \ifdim\wd\nmFigBox>\linewidth
    \setbox\nmFigBox=\hbox{\resizebox{\linewidth}{!}{\usebox{\nmFigBox}}}%
  \fi
  \@tempdimc=\ht\nmFigBox \advance\@tempdimc by \dp\nmFigBox
  \ifdim\@tempdimc>#1\relax
    \setbox\nmFigBox=\hbox{\resizebox{!}{#1}{\usebox{\nmFigBox}}}%
  \fi}
\newcommand{\nmfigureh}[2]{%
  \begingroup
  \begin{lrbox}{\nmFigBox}\input{#2}\end{lrbox}%
  \nmFitBox{#1}%
  \begin{center}\usebox{\nmFigBox}\end{center}%
  \endgroup}
\newcommand{\nmfigure}[1]{\nmfigureh{\textheight}{#1}}
% \lkimageh{4cm}{file.png}: a raster or PDF image with the same rule (a clip's still, say)
\newcommand{\lkimageh}[2]{%
  \begingroup
  \begin{lrbox}{\nmFigBox}\includegraphics{#2}\end{lrbox}%
  \nmFitBox{#1}%
  \begin{center}\usebox{\nmFigBox}\end{center}%
  \endgroup}
\makeatother
% a labelled placeholder, so that a deck with a missing figure still builds
\newcommand{\lkplaceholder}[2]{%
  \begin{center}%
    \begin{tikzpicture}%
      \draw[nmCaveat,dashed,thin] (0,0) rectangle (0.8\linewidth,#1);%
      \node[nmCaveat,font=\sffamily\scriptsize] at (0.4\linewidth,0.5*#1) {#2};%
    \end{tikzpicture}%
  \end{center}}
% Figures lifted from papers keep their white background; on a warm-paper slide that reads as a white card.
% Compositing the image with the multiply blend mode turns its white into the slide background and leaves ink
% and colours in place; on a white slide it changes nothing. Set without a transparency group on purpose.
\newsavebox{\nmFigureBox}
\newenvironment{paperfigure}{\begin{lrbox}{\nmFigureBox}}{\end{lrbox}%
  \begin{tikzpicture}[blend mode=multiply]\node[inner sep=0pt,outer sep=0pt]{\usebox{\nmFigureBox}};\end{tikzpicture}}
% Short caveat callouts retain charcoal body text; the rule and label encode role.
\newcommand{\caveatnote}[1]{%
  \par\vspace{2mm}\noindent
  \begin{minipage}[t]{.01\textwidth}\vspace{0pt}{\color{nmCaveat}\rule{1pt}{13mm}}\end{minipage}\hspace{3mm}%
  \begin{minipage}[t]{.91\textwidth}\vspace{0pt}\raggedright
  {\color{nmCaveat}\sffamily\mdseries Caveat}\par\smallskip
  {\color{nmInk}#1}\end{minipage}\par}

% ---------------------------------------------------------------- citations and references
% Every in-text citation is set in italics (rule of 27 October 2026), and smaller with citation-size: small; the
% dual-output filter wraps the citations, so nothing is needed here.
\setbeamertemplate{bibliography item}[text]
% Pandoc 3.1.7 and later emit the references as Beamer \bibitem entries, which take their colours from these
% templates (the default is the structure colour); keep them in ink.
\setbeamercolor{bibliography entry author}{fg=nmInk}
\setbeamercolor{bibliography entry title}{fg=nmInk}
\setbeamercolor{bibliography entry location}{fg=nmInk}
\setbeamercolor{bibliography entry note}{fg=nmInk}
\setbeamerfont{bibliography entry author}{size=\footnotesize}
\setbeamerfont{bibliography entry title}{size=\footnotesize}
\setbeamerfont{bibliography entry location}{size=\footnotesize}
\setbeamerfont{bibliography entry note}{size=\footnotesize}
% References in footnote size with half a line between entries, so that a frame holds five or six entries rather
% than two or three. Pandoc defines CSLReferences before this header.
\makeatletter
\@ifundefined{CSLReferences}{}{%
  \renewenvironment{CSLReferences}[2]{\footnotesize\begin{list}{}{%
    \setlength{\itemindent}{0pt}\setlength{\leftmargin}{0pt}\setlength{\parsep}{0pt}%
    \ifodd #1\setlength{\leftmargin}{\cslhangindent}\setlength{\itemindent}{-1\cslhangindent}\fi
    \setlength{\itemsep}{0.45\baselineskip}}}{\end{list}}}
\makeatother
\hypersetup{colorlinks=true,urlcolor=nmInput,linkcolor=nmInput,citecolor=nmState}
% speaker notes never reach the study PDF (the filter removes them; this covers raw \note{} too)
\setbeameroption{hide notes}
% the deck's own metadata (short title, eyebrow, citation size), written by the filter per deck
\InputIfFileExists{_lecturekit/\jobname-deck.tex}{}{}
